\documentclass[11pt]{article}

% Math
\usepackage{amsmath, amssymb, amsfonts, mathtools}

% Page layout
\usepackage[margin=1in]{geometry}

% Optional useful packages
\usepackage{physics}      % \bra, \ket, \dv, \pdv, etc.
\usepackage{bm}           % bold math symbols
\usepackage{enumitem}     % better lists
\usepackage{hyperref}     % links
\hypersetup{
    colorlinks=true,       % false: boxed links; true: colored links
    linkcolor=blue,        % color of internal links (sections, pages, etc.)
    citecolor=green,       % color of citation links to bibliography
    % filecolor=magenta,     % color of file links
    urlcolor=blue,         % color of external website links
    % allcolors=blue       % uncomment this if you want ALL links to use one color
}

\title{DB Tutorial}
\author{Reza Asakereh}
\date{Last Update: \today}

\begin{document}

\maketitle

\section{Findings}
\subsection{Move toward ground state}
I had an interesting observation regarding DB-QITE:
The purpose of this algorithm is that when applied to state $|\psi_{old}\rangle$ we get $|\psi_{new}\rangle$ that is closer to the ground state of $H$ than before.

To measure how close a state ($|\psi\rangle$) is to the ground state ($|v_0\rangle$) we can use the size of dot product: $| \langle v_0 | \psi \rangle |$. So the purpose of DB-QITE is to ensure
$$| \langle v_0 | \psi_{new} \rangle | > | \langle v_0 | \psi_{old} \rangle |$$

My observation is that if we show $| \psi_{old} \rangle = (a_0, a_1, \dots, a_n)$ in the basis of $H$ eigenvectors, we will have:
$$
| \langle v_0 | \psi_{new} \rangle |^2 = | \langle v_0 | \psi_{old} \rangle |^2 + t(\sum_{i=0}^n{2(s_i - s_0)|a_i|^2|a_0|^2}) + O(t^2)
$$
where $s_i$ is the $i$'th smallest eigenvalue of $H$ and $t$ is time.
\\

It is obvious that the value added to $| \langle v_0 | \psi_{old} \rangle |^2$ is always positive except for when $|a_0| = 0$ or $1$ therefore we always get closer to the ground state.
Don't worry about the two edge cases. They mean $| \psi_{old} \rangle$ has no ground state component or is already ground state itself, so for the former convergence never happens using ITE algorithms and for the latter we have already converged.
\\


\textbf{Intuition:}
We can see that change in $| \langle v_0 | \psi \rangle |^2$ is proportional to $\sum_{i=0}^n{(s_i - s_0)|a_i|^2}$. This sum actually has an interesting physical interpretation. It is equal to $\langle \psi | H | \psi \rangle - s_0$ which means how far the average energy of the state is from the ground state. This says that for the same starting $| \langle v_0 | \psi \rangle |$, the states with higher average energy, converge faster towards the ground state.

Details and proof can be found in \href{https://colab.research.google.com/drive/1dsnv5_XMr6TzfRn7-RQuOCWzfpAqkZ1y?usp=sharing}{here}

\end{document}